\subsection{诺特环上的有限生成模情形}

定理 1. 设 A 是诺特环，M 是有限生成的 A-模。存在 M 的一个合成列 \((\mathbf{M}_i)_{0\leqslant i\leqslant n}\)，使得对于 \(0\leqslant i\leqslant n - 1\)，\(\mathbf{M}_i / \mathbf{M}_{i + 1}\) 同构于 \(\mathbf{A} / \mathfrak{p}_i\)，其中 \(\mathfrak{p}_i\) 是 A 的素理想。

令 \(\mathfrak{E}\) 为 M 的子模所成的集合，这些子模具有满足陈述的合成列。由于 \(\mathfrak{E}\) 非空（因为 \{0\} 属于 \(\mathfrak{E}\)），且 M 是诺特模，\(\mathfrak{E}\) 有极大元 N。若 M ≠ N，则 M/N ≠ 0，从而 \(\operatorname{Ass}(M/N) \neq \varnothing\)（第 1 节，命题 2 的推论 1）；因此 M/N 含有一个子模 N'/N，同构于形如 A/p 的 A-模，其中 p 是素理想；于是按定义 N' \(\mathfrak{E}\)，这与 N 的极大性矛盾。因此必有 N = M。

定理 2. 设 M 是诺特环 A 上的有限生成模，\((\mathbf{M}_{i})_{0\leqslant i\leqslant n}\) 是 M 的一个合成列，使得对于 \(0\leqslant i\leqslant n-1\)，\(M_{i}/M_{i+1}\) 同构于 \(A/p_{i}\)，其中 \(p_{i}\) 是 A 的素理想。则

\[
\operatorname{Ass} (\mathbf {M}) \subset \left\{\mathfrak {p} _ {0}, \dots , \mathfrak {p} _ {n - 1} \right\} \subset \operatorname{Supp} (\mathbf {M});\tag{4}
\]

这三个集合的极小元相同，并且与包含 Ann(M) 的素理想集合的极小元重合。

包含关系 \(\operatorname{Ass}(\mathbf{M}) \subset \{\mathfrak{p}_0, \ldots, \mathfrak{p}_{n-1}\}\) 立即由第 1 节的命题 1 和命题 3 得出。对于 \(0 \leqslant i \leqslant n - 1\)，

\[
\mathfrak {p} _ {i} \in \operatorname{Supp} (\mathrm{A} / \mathfrak {p} _ {i}) = \operatorname{Supp} (\mathrm{M} _ {i} / \mathrm{M} _ {i + 1})
\]

（第二章，第 4 节，第 4 节，例），从而 \(\mathfrak{p}_i \in \operatorname{Supp}(\mathbf{M}_i) \subset \operatorname{Supp}(\mathbf{M})\)（第二章，第 4 节，第 4 节，命题 16），这就说明包含关系

\[
\left\{\mathfrak {p} _ {0}, \dots , \mathfrak {p} _ {n - 1} \right\} \subset \operatorname{Supp} (\mathbf {M}).
\]

第 1 节第 3 节命题 7 的推论 1 表明，\(\operatorname{Ass}(\mathbf{M})\) 和 \(\operatorname{Supp}(\mathbf{M})\) 具有相同的极小元，而 (4) 表明这些极小元正是 \(\{\mathfrak{p}_0, \ldots, \mathfrak{p}_{n-1}\}\) 的极小元。最后一个断言由第二章，第 4 节，第 4 节，命题 17 得出。

推论。若 M 是诺特环上的有限生成模，则 Ass(M) 是有限集。

在定理 2 的条件下，集合 \(\{\mathfrak{p}_0,\ldots ,\mathfrak{p}_{n - 1}\}\) 不一定由 M 唯一确定；特别地，它可能不同于 Ass(M)（习题 6）。

命题 8. 设 A 是诺特环，a 是 A 的一个理想，M 是有限生成的 A-模。下列条件等价：

\begin{enumerate}
\def\labelenumi{(\alph{enumi})}
\item
存在一个元素 \(x \neq 0\)，它属于 \(M\)，并且满足 \(ax = 0\)。
\item
对于所有 \(a \in \mathfrak{a}\)，存在一个元素 \(x \neq 0\)，它属于 \(\mathbf{M}\)，且满足 \(ax = 0\)；
\item
  存在 \(\mathfrak{p} \in \operatorname{Ass}(\mathbf{M})\)，使得 \(\mathfrak{a} \subset \mathfrak{p}\)。
\end{enumerate}

显然 (a) 蕴含 (b)。根据第 1 节命题 2 的推论 2，条件 (b) 意味着理想 a 包含于与 M 相伴的素理想的并中；又因为 Ass(M) 是有限的（第二章，第 1 节，第 1 节，命题 2），故它包含于其中一个素理想，从而 (b) 蕴含 (c)。最后，若存在 p ∈ Ass(M) 使得 a ⊂ p，则 p 是 M 中某个元素 x ≠ 0 的湮灭子（第 1 节，定义 1），且 ax = 0；因此 (c) 蕴含 (a)。

命题 9. 设 A 是诺特环，a 是 A 的一个理想，M 是有限生成的 A-模。存在整数 n \textgreater{} 0 使得 \(a^{n}M = 0\) 的充要条件是：a 包含于与 M 相伴的素理想的交中。

这个交也是 Supp(M) 的极小元的交（第 3 节，命题 7 的推论 1）；说 a 包含于这个交中，等价于按第二章，第 4 节的记号有 V(a) ⊃ Supp(M)；结论于是由第二章，第 4 节，第 4 节，命题 17 的推论 2 得出。

定义 2. 给定 A-模 M，M 的一个自同态 u 称为殆幂零的，如果对于所有 \(x \in M\)，存在整数 \(n(x) > 0\)，使得 \(u^{n(x)}(x) = 0\)。

如果 M 是有限生成的，则每个殆幂零自同态都是幂零的。

推论。设 A 是诺特环，M 是 A-模，a 是 A 的一个元素。要使同态 \(a_{\mathbf{M}}: x \mapsto ax\) 作用于 \(\mathbf{M}\) 而成为殆幂零，充要条件是 a 属于 \(\operatorname{Ass}(\mathbf{M})\) 的每个理想。

陈述中的条件等价于说，对于所有 \(x \in M\)，存在 \(n(x) > 0\)，使得 \((\mathbf{A}a)^{n(x)}(\mathbf{A}x) = 0\)；根据命题 9，这也意味着 a 属于与子模 Ax 相伴的所有素理想；推论于是由以下事实得出：当 x 遍历 M 时，\(\operatorname{Ass}(M)\) 是所有 \(\operatorname{Ass}(Ax)\) 的并（第 1 节，公式 (1)）。

命题 10. 设 A 是诺特环，E 是有限生成的 A-模，F 是 A-模。则

\[
\operatorname{Ass} \left(\operatorname{Hom} _ {\mathbf {A}} (\mathrm{E}, \mathrm{F})\right) = \operatorname{Ass} (\mathrm{F}) \cap \operatorname{Supp} (\mathrm{E}).\tag{5}
\]

根据假设，E 同构于形如 \(A^{n}/R\) 的 A-模，因此 \(\operatorname{Hom}_{\mathbf{A}}(\mathbf{E}, \mathbf{F})\) 同构于 \(\operatorname{Hom}_{\mathbf{A}}(\mathbf{A}^{n}, \mathbf{F})\) 的一个子模，而后者同构于 \(F^{n}\)；现在 \(\operatorname{Ass}(\mathbf{F}^{n}) = \operatorname{Ass}(\mathbf{F})\)（第 1 节，命题 3 的推论 1）；所以 \(\operatorname{Ass}(\operatorname{Hom}_{\mathbf{A}}(\mathbf{E}, \mathbf{F})) \subset \operatorname{Ass}(\mathbf{F})\)。另一方面，

\[
\operatorname{Supp} (\operatorname{Hom} _ {\mathbf {A}} (\mathbf {E}, \mathbf {F})) \subset \operatorname{Supp} (\mathbf {E}):
\]

对于 A 的每个素理想 p，\(\operatorname{Hom}_{\mathbf{A}_{p}}(\mathbf{E}_{p}, \mathbf{F}_{p})\) 同构于 \((\operatorname{Hom}_{\mathbf{A}}(\mathbf{E}, \mathbf{F}))_{p}\)（第二章，第 2 节，第 7 节，命题 19），我们的断言立即得出；于是由定理 2 得

\[
\operatorname{Ass} \left(\operatorname{Hom} _ {\mathbf {A}} (\mathrm{E}, \mathrm{F})\right) \subset \operatorname{Supp} (\mathrm{E}).
\]

反之，设 p 是属于 \(\operatorname{Ass}(\mathbf{F}) \cap \operatorname{Supp}(\mathbf{E})\) 的 A 的素理想。根据定义，F 含有一个同构于 A/p 的子模。~另一方面，由于 E 有限生成且 \(E_{p} \neq 0\)，我们知道存在从 E 到 A/p 的非零同态 \(w \neq 0\)（第二章，第 4 节，第 4 节，命题 20）。由于存在从 A/p 到 F 的单射同态 j，\(j \circ w \in \operatorname{Hom}(\mathbf{E}, \mathbf{F})\) 且 \(j \circ w \neq 0\)。另一方面，对于某个 \(a \in A\)，关系 aw = 0 等价于 \(a \in p\)；A/p 中每个元素 \(\neq 0\) 的湮灭子都是 p；于是当然有 \(p \in \operatorname{Ass}(\operatorname{Hom}_{\mathbf{A}}(\mathbf{E}, \mathbf{F}))\)。

